% GENERATED by scripts/build-lecture-slides.py; do not edit. Portable profile: XeTeX, theme and fonts from the core.
\makeatletter\IfFileExists{core/beamer/beamerthememx.sty}{\def\mxcorepath{core/}\def\input@path{{core/beamer/}{core/}}}{\def\mxcorepath{../core/}\def\input@path{{../core/beamer/}{../core/}}}\makeatother
\documentclass[aspectratio=169,10pt,t]{beamer}
\usetheme{mx}
\usepackage{graphicx,booktabs}
\usepackage[french]{babel}
\setlength{\parskip}{.45em}
\hypersetup{pdfauthor={},pdfsubject={Lecture slides},pdfkeywords={},pdftitle={Applications affines et isométries}}
\providecommand{\N}{\mathbb{N}}
\providecommand{\Z}{\mathbb{Z}}
\providecommand{\D}{\mathbb{D}}
\providecommand{\Q}{\mathbb{Q}}
\providecommand{\R}{\mathbb{R}}
\providecommand{\C}{\mathbb{C}}
\providecommand{\vect}[1]{\overrightarrow{#1}}
\providecommand{\vi}{\vec{\imath}}
\providecommand{\vj}{\vec{\jmath}}
\providecommand{\vk}{\vec{k}}
\providecommand{\norm}[1]{\left\lVert #1\right\rVert}
\providecommand{\vc}[2]{\begin{pmatrix}#1\\#2\end{pmatrix}}
\providecommand{\deter}[2]{\det\left(#1,#2\right)}
\providecommand{\abs}[1]{\left\lvert #1\right\rvert}
\providecommand{\dd}{\,\mathrm{d}}
\providecommand{\e}{\mathrm{e}}
\providecommand{\Card}{\operatorname{Card}}
\providecommand{\E}{\operatorname{E}}
\providecommand{\V}{\operatorname{V}}
\providecommand{\degre}{^{\circ}}
\providecommand{\un}[1]{\,\mathrm{#1}}
\providecommand{\sci}[2]{#1\times10^{#2}}
\providecommand{\pgcd}{\operatorname{PGCD}}
\providecommand{\ppcm}{\operatorname{PPCM}}
\title{Applications affines et isométries}
\subtitle{Cours et exemples}
\date{}
\author{}
\begin{document}
\mxtitleframe
\begin{frame}{Introduction}
Dans ce chapitre, \(\mathcal{P}\) désigne le plan affine et \(\mathcal{E}\) l'espace affine.
\end{frame}
\begin{frame}{Définitions et exemples}
Une application \(f : \mathcal{P} \rightarrow \mathcal{P}\) est une \emph{application affine}
                  si elle conserve le barycentre.
                  
                  Si \(G = \mathrm{Bar}\left\lbrace (A,\alpha) ; (B,\beta) \right\rbrace\).
                  En posant \(G' = f(G)\), \(A' = f(A)\) et \(B' = f(B)\), on a 
                  \(G' = \mathrm{Bar}\left\lbrace (A',\alpha) ; (B',\beta) \right\rbrace\).
\end{frame}
\begin{frame}{Définitions et exemples}
\begin{enumerate}
\setcounter{enumi}{0}\item Considérons la translation \(t_{\overrightarrow{u}}\) de vecteur \(\overrightarrow{u}\).
                              En posant \(t_{\overrightarrow{u}} (M) = M'\), on a \(\overrightarrow{MM'} = \overrightarrow{u}\).
                              
                              Soit \(G = \mathrm{Bar}\left\lbrace (A,\alpha) ; (B,\beta) \right\rbrace\).
                              Soient
                              \(G' = t_{\overrightarrow{u}} (G)\),
                              \(A' = t_{\overrightarrow{u}} (A)\) et
                              \(B' = t_{\overrightarrow{u}} (B)\)
                              .
                              On a alors
                              \[\begin{aligned}
                              \alpha \overrightarrow{G'A'} + \beta \overrightarrow{G'B'}
                                    & =  \alpha (\overrightarrow{G'G} + \overrightarrow{GA} + \overrightarrow{AA'})
                                        +\beta (\overrightarrow{G'G} + \overrightarrow{GB} + \overrightarrow{BB'})      \\
                                    & =  \alpha (-\overrightarrow{u} + \overrightarrow{GA} + \overrightarrow{u})
                                        +\beta (-\overrightarrow{u} + \overrightarrow{GB} + \overrightarrow{u})      \\
                                    & = \alpha \overrightarrow{GA} + \beta \overrightarrow{GB}      \\
                                    & = \overrightarrow{0}.
                              \end{aligned}\]
                              Donc \(G' = \mathrm{Bar}\left\lbrace (A';\alpha);(B';\beta) \right\rbrace\) et ainsi, \(t_{\overrightarrow{u}}\) est une application affine.

\end{enumerate}
\end{frame}
\begin{frame}{Définitions et exemples}
Considérons l'homothétie \(h\) de centre \(I\) et de rapport \(k\).
                  En posant \(h(M) = M'\), on a \(\overrightarrow{IM'} = k \overrightarrow{IM}\).
                  
                  On veut montrer que \(h\) est une application affine.
                  Considérons le point
                  \[
                  G = \mathrm{Bar}\left\lbrace (A,\alpha);(B,\beta) \right\rbrace.
                  \]
                  Posons \(G' = h(G)\), \(A' = h(A)\) et \(B' = h(B)\).
                  On a alors 
                  \[\begin{aligned}
                        \alpha \overrightarrow{G'A'} + \beta \overrightarrow{G'B'}
                              & =  \alpha (\overrightarrow{G'I} + \overrightarrow{IA'})
                                  +\beta (\overrightarrow{G'I} + \overrightarrow{IB'})      \\
                              & = \alpha (k \overrightarrow{GI} + k \overrightarrow{IA}) 
                                  +\beta (k \overrightarrow{GI} + k \overrightarrow{IB})      \\
                              & = k \left[ \alpha \overrightarrow{GA} + \beta \overrightarrow{GB} \right]       
                                = k \overrightarrow{0}      
                                = \overrightarrow{0}.
                  \end{aligned}\]
\end{frame}
\begin{frame}{Définitions et exemples}
Considérons la réflexion d'axe \((\mathcal{D})\).
                  On veut montrer que la réflexion est une application affine.

                  Soit \(G = \mathrm{Bar}\left\lbrace (A,\alpha);(B,\beta) \right\rbrace\).
                  Alors, \(A\), \(G\) et \(B\) sont alignés.
                  \[\begin{aligned}
                        \alpha \overrightarrow{GA} + \beta \overrightarrow{GB} & = \overrightarrow{0} \\
                        \alpha \overrightarrow{GA} + \beta \overrightarrow{GA} + \beta \overrightarrow{AB} & = \overrightarrow{0} \\
                        (\alpha+\beta) \overrightarrow{GA} + \beta \overrightarrow{AB} & = \overrightarrow{0}
                  \end{aligned}\]
                  et donc \(\overrightarrow{AG} = \dfrac{\beta}{\alpha + \beta} \overrightarrow{AB}\).
                  Comme
                  \(\overline{AG} = \dfrac{\beta}{\alpha + \beta} \overline{AB}\), on a 
                  \(\dfrac{\overline{AG}}{\overline{AB}} = \dfrac{\beta}{\alpha + \beta}\).
                  
                  De plus, \((AA') \parallel (BB') \parallel (GG')\),
                  donc 
                  \(\dfrac{\overline{AG}}{\overline{AB}} = \dfrac{\overline{A'G'}}{\overline{A'B'}}\).
                  D'où
                  \(\dfrac{\overline{A'G'}}{\overline{A'B'}} = \dfrac{\beta}{\alpha+\beta}\)
                  et, par conséquent,
                  \(\overrightarrow{A'G'} = \dfrac{\beta}{\alpha+\beta} \overrightarrow{A'B'}\).
                  
                  Ainsi, \(G' = \mathrm{Bar}\left\lbrace (A',\alpha);(B',\beta) \right\rbrace\).
\end{frame}
\begin{frame}{Définitions et exemples}
Considérons la rotation \(r\) de centre \(I\) et d'angle \(\theta\).
                  Elle a alors pour expression complexe \(z' = az + b\) où \(a = e^{i\theta}\), donc \(|a| = 1\).
                  
                  Soit \(G = \mathrm{Bar}\left\lbrace (A,\alpha);(B,\beta) \right\rbrace\). 
                  Posons
                        \[
                        G' = r(G)
                              \qquad
                        A' = r(A)
                              \qquad
                        B' = r(B).
                        \]
                  L'affixe de \(\alpha \overrightarrow{G'A'} + \beta \overrightarrow{G'B'}\) 
                  est \(\alpha (z_A' - z_G') + \beta (z_B' - z_G')\),
                  ce qui est égal à 
                  \[
                        \alpha (a z_A + b - a z_G - b) + \beta (a z_B + b - a z_G - b) =
                                                      a \left[ \alpha (z_A - z_G) + \beta (z_B - z_G) \right] \]
                  qui n'est autre que l'affixe de 
                  \(a ( \underbrace{\alpha \overrightarrow{GA} + \beta \overrightarrow{GB}}_{\overrightarrow{0}} )\).
                  
                  Ainsi, \(G'= \mathrm{Bar}\left\lbrace (A',\alpha);(B',\beta) \right\rbrace\).
\end{frame}
\begin{frame}{Définitions et exemples}
La composition de deux applications affines est une application affine.
             
            C'est-à-dire que si \(f\) et \(g\) sont deux applications affines, alors \(f\circ g\) est aussi une application affine.
\end{frame}
\begin{frame}{Définitions et exemples}
Soit \(f : \mathcal{P} \rightarrow \mathcal{P} , M \mapsto M'\) une application affine.
            L'application \(\Phi : \overrightarrow{AB} \mapsto \Phi(\overrightarrow{AB}) = \overrightarrow{A'B'}\)
            est appelée application linéaire associée à l'application affine \(f\).
\end{frame}
\begin{frame}{Définitions et exemples}
Trouvons l'application linéaire associée à une translation \(t_{\overrightarrow{u}}\).

Appelons \(T\) l'application linéaire associée à \(t_{\overrightarrow{u}}\).
                  On a alors \(T : \overrightarrow{AB} \mapsto \overrightarrow{A'B'} = \overrightarrow{AB}\).
                  Ainsi, \(T = \mathrm{Id}\).
\end{frame}
\begin{frame}{Définitions et exemples}
Cherchons l'application linéaire associée à une homothétie \(h(I,k)\).

                  Soit \(H\) l'application linéaire associée à \(h(I,k)\).
                  On a \(H : \overrightarrow{AB} \mapsto \overrightarrow{A'B'} = k \overrightarrow{AB}\).
                  
                  L'application \(H\) est appelée une \emph{homothétie vectorielle}.
\end{frame}
\begin{frame}{Définitions et exemples}
Soit \(f\) l'application affine qui à tout point \(M(x,y)\) associe le point \(M'(x',y')\).
                  On appelle \emph{expression analytique} de \(f\) la donnée de \((x',y')\) en fonction de \((x,y)\).
                  Elle est souvent donnée sous la forme : 
                  \[
                        \left\lbrace
                              \begin{array}{l}
                                    x' = a_1 x + a_2 y + a_3 \\
                                    y' = b_1 x + b_2 y + b_3 
                              \end{array}
                        \right.
                        
                  \]
                  où \(a_1\), \(a_2\), \(a_3\), \(b_1\), \(b_2\) et \(b_3\) sont des réels.
\end{frame}
\begin{frame}{Définitions et exemples}
Une application affine \(f\) du plan \(\mathcal{P}\) est entièrement définie par les données 
                  des transformés de trois points non-alignés \(A\), \(B\) et \(C\).
\end{frame}
\begin{frame}{Définitions et exemples}
[.64]
			                  Soient \(A\), \(B\) et \(C\) trois points non-alignés du plan.
			                  Alors \((A ; \overrightarrow{AB} , \overrightarrow{AC})\) forme un repère.
      
                              Pour tout point \(M\) du plan, il existe \((x,y)\) tel que 
                              \(\overrightarrow{AM} = x \overrightarrow{AB} + y \overrightarrow{AC}\).
			                  Comme \(f\) est une application affine,
			                  \(\overrightarrow{A'M'} = x \overrightarrow{A'B'} + y \overrightarrow{A'C'}\),
			                  et donc
			            
			                  !

                        \[
                        \begin{pmatrix} x' - x_A' \\ y' - y_A' \end{pmatrix}
                        =
                        x \begin{pmatrix} x_B' - x_A' \\ y_B' - y_A' \end{pmatrix}
                        +
                        y \begin{pmatrix} x_C' - x_A' \\ y_C' - y_A' \end{pmatrix}.
                        \]
\end{frame}
\begin{frame}{Définitions et exemples (2)}
Ce qui détermine \((x',y')\).
\end{frame}
\begin{frame}{Isométrie affine}
Une transformation \(f\) est une isométrie si elle conserve la distance, autrement dit :
            si 
            \( \left\{\begin{array}{l}
                  f(M) = M' \\
                  f(N) = N'
            \end{array}\right.      
            \)
            alors 
            \(M'N' = MN\).
            
            Nous admettons que toute \emph{isométrie} est une \emph{application affine bijective}.
\end{frame}
\begin{frame}{Isométrie affine}
\begin{itemize}
\item \textbf{Translation :} 

                        Soit \(\overrightarrow{u}\) un vecteur et soient \(M\) et \(N\) deux points du plan.
                        Notons par \(t_{\overrightarrow{u}}\) la translation de vecteur \(\overrightarrow{u}\). 
                        On a alors 
                        \[\begin{array}{rcl}
                              t_{\overrightarrow{u}} (M) = M' & \Rightarrow & \overrightarrow{MM'} = \overrightarrow{u} \\
                              t_{\overrightarrow{u}} (N) = N' & \Rightarrow & \overrightarrow{NN'} = \overrightarrow{u} \\
                        \end{array}\]
                        Ainsi,
                        \[
                              \overrightarrow{M'N'} = \underbrace{\overrightarrow{M'M}}_{-\overrightarrow{u}}
                                                      + \overrightarrow{MN}
                                                      + \underbrace{\overrightarrow{NN'}}_{\overrightarrow{u}}
                                       \Rightarrow
                                                      \overrightarrow{M'N'} = \overrightarrow{MN}
                        \]
                        et donc \(M'N' = MN\).

\end{itemize}
\end{frame}
\begin{frame}{Isométrie affine (2)}
\begin{itemize}
\item \textbf{Rotation :} 

                        Considérons une rotation \(r\) de centre \(I\).
                        Si \(M' = r(M)\) et \(N' = r(N)\) alors les triangles \(IMN\) et \(IM'N'\) sont superposables.
                        Ainsi, \(MN = M'N'\).

\end{itemize}
\end{frame}
\begin{frame}{Isométrie affine (3)}
\begin{itemize}
\item \textbf{Reflexion :} 

                        Considérons la symétrie orthogonale d'axe \((\mathcal{D})\).
                        
                        Les trapèzes \(ABNM\) et \(ABN'M'\) sont superposables (symétriques l'un de l'autre par rapport à \((\mathcal{D})\)), donc \(M'N' = MN\).

\end{itemize}
\end{frame}
\begin{frame}{Isométrie affine (4)}
\begin{itemize}
\item \textbf{Homothétie :} 

                        Considérons l'homothétie de centre \(I\) et de rapport \(k\).

                        Comme \(\overrightarrow{IM'} = k \overrightarrow{IM}\) et \(\overrightarrow{IN'} = k \overrightarrow{IN}\), 
                        on a \(\dfrac{IM'}{IM} = |k|\) et \(\dfrac{IN'}{IN} = |k|\).
                        
                        Comme \(\dfrac{IM'}{IM} = \dfrac{IN'}{IN}\), 
                        d'après la réciproque de la propriété de Thalès,
                        \((MN) \parallel (M'N')\).
                        On a alors \(\dfrac{IM'}{IM} = \dfrac{IN'}{IN} = \dfrac{M'N'}{MN} = |k|\)
                        (plus directement : \(\overrightarrow{M'N'} = \overrightarrow{IN'} - \overrightarrow{IM'} = k\overrightarrow{MN}\)).
                        
                        Une homothétie de rapport \(k\) est une isométrie si et seulement si \(|k| = 1\).

\end{itemize}
\end{frame}
\begin{frame}{Isométrie affine}
\begin{itemize}
\item La composition de deux isométries est une isométrie.
                  
\item La réciproque d'une isométrie est une isométrie.
                  
\item 
\begin{itemize}

                              \item L'image d'une droite par une isométrie est une droite.
                              
\item Toute isométrie conserve l'alignement.

\end{itemize}

\end{itemize}
\end{frame}
\begin{frame}{Isométrie affine (2)}
\begin{itemize}
\item Toute isométrie transforme deux droites parallèles en deux droites parallèles :
                        on dit qu'une isométrie \emph{conserve le parallélisme}.
                  
\item Toute isométrie transforme deux droites perpendiculaires en deux droites perpendiculaires :
                        on dit qu'une isométrie \emph{conserve l'orthogonalité}.

\end{itemize}
\end{frame}
\begin{frame}{Isométrie affine (3)}
\begin{itemize}
\item Toute isométrie transforme un cercle en un cercle : 
                        on dit qu'une isométrie \emph{conserve la cocyclicité}.
                  
\item Toute isométrie transforme 2 cercles tangents 
                        (respectivement un cercle et une droite tangente)
                        en 2 cercles tangents
                        (respectivement en un cercle et une droite tangente) :
                        on dit qu'une isométrie \emph{conserve le contact}.

\end{itemize}
\end{frame}
\begin{frame}{Reconnaissance des isométries}
Soit \(f\) une application affine d'expression analytique 
            \[
                  \left\lbrace
                        \begin{array}{l}
                              x' = a_1 x + a_2 y + a_3 \\
                              y' = b_1 x + b_2 y + b_3 \;.
                        \end{array}
                  \right.
            \]

Soit 
            \( A = \begin{pmatrix} a_1 & a_2 \\ b_1 & b_2 \end{pmatrix} \)
            et soit \(\Delta = a_1 \cdot b_2 - b_1 \cdot a_2\)
            .
\end{frame}
\begin{frame}{Reconnaissance des isométries}
Le repère étant orthonormé, l'application \(f\) est une isométrie si et seulement si \(A\) est de l'un des types : 
                  \[\begin{aligned}
                        \begin{pmatrix}
                              a & -b \\ b & a 
                        \end{pmatrix}
                              \qquad \text{ ou } \qquad
                        \begin{pmatrix}
                              a & b \\ b & -a
                        \end{pmatrix}
                  \end{aligned}\]
                  où \(a^2 + b^2 = 1\).
                  
                  De plus, \(f\) est une isométrie positive (ou un déplacement) si \(\Delta = 1\)
                  et elle est une isométrie négative (ou un antidéplacement) si \(\Delta = -1\).
\end{frame}
\begin{frame}{Reconnaissance des isométries}
Soit \(f\) une application affine.

\begin{itemize}

                  \item L'application \(f\) est une isométrie positive si son expression complexe est de la forme
                        \(z' = a z + b\) où \(|a| = 1\).
                  
\item L'application \(f\) est une isométrie négative si son expression complexe est de la forme
                        \(z' = a \overline{z} + b\) où \(|a| = 1\).

\end{itemize}
\end{frame}
\begin{frame}{Reconnaissance des isométries}
Soit \(f\) une isométrie.
		      table[h!]
\end{frame}
\begin{frame}{Reconnaissance des isométries (2)}
\begin{center}\begin{tabular}{lll}
\textbf{Points invariants} & \textbf{déplacement} & \textbf{antidéplacement} \\ 
Le plan \(\mathcal{P}\) tout entier & \(\mathrm{Id}\) &  \\ 
Une droite \((\mathcal{D})\) &  & \(S_{(\mathcal{D})}\) \\ 
Un point \(\Omega\) & \(r(\Omega,\theta),\ \theta\neq 0\) &  \\ 
Aucun point invariant & \(t_{\overrightarrow{u}},\ \overrightarrow{u}\neq\overrightarrow{0}\) &  
                                            Composée d'une symétrie orthogonale d'axe \((\mathcal{D})\) et 
                                            d'une translation de vecteur \(\overrightarrow{u}\) non nul de 
                                            direction parallèle à \((\mathcal{D})\).
                                                 [0] \\ 
\end{tabular}\end{center}
\end{frame}
\begin{frame}{Reconnaissance des isométries (3)}
Classification des isométries suivant l'ensemble des points invariants.
            
            table
\end{frame}
\begin{frame}{Classification des isométries}
Soit \(\left(\overrightarrow{OA},\overrightarrow{OB}\right)\) un angle orienté.

\begin{itemize}

            \item La translation ne change pas l'orientation des angles.

\item La rotation ne change pas l'orientation des angles.

\item La réflexion change l'orientation des angles.

\end{itemize}

Toute isométrie qui \emph{conserve l'orientation des angles} est appelée \emph{déplacement}.
\end{frame}
\begin{frame}{Classification des isométries (2)}
Toute isométrie qui \emph{\textbf{ne} conserve \textbf{pas} l'orientation des angles} est appelée

\emph{\textbf{anti}déplacement}.
\end{frame}
\begin{frame}{Classification des isométries}
\begin{itemize}
\item Toute isométrie est soit un déplacement, soit un antidéplacement.
                  
\item La réciproque d'un déplacement est un déplacement.
                  
\item La réciproque d'un antidéplacement est un antidéplacement.
                  
\item La composition de deux déplacements est un déplacement.
                  
\item La composition de deux antidéplacements est un déplacement.

\end{itemize}
\end{frame}
\begin{frame}{Classification des isométries (2)}
\begin{itemize}
\item La composition d'un déplacement et d'un antidéplacement,
                        dans n'importe quel ordre, est un antidéplacement.

\end{itemize}
\end{frame}
\begin{frame}{Classification des isométries}
Tout déplacement a pour expression complexe une équation de la forme \(z' = az + b\) avec \(|a| = 1\).
	             
	            0.96!\(\xymatrix@C=.05\actualtextwidth{
                            &
                                   &  
                                          & *+[F]{ \begin{array}{c} f:z' = az + b \\ |a| = 1 \end{array} }
                                            \ar[dr]
                                            \ar[dll] 
                                                 &
                                                        \\
                            & *+[F]{a=1} 
                              \ar[dr]
                              \ar[dl]
                                   &  
                                          &
                                                 & *+[F]{a \neq 1 }
                                                   \ar[dd]
                                                        \\
                     *+[F]{b=0} 
                     \ar[d]
                            &            
                                   & *+[F]{b\neq 0}
                                     \ar[d]
                                          &
                                                 &
                                                        \\ 
                     *+[F]{f=\mathrm{Id}}
                            & 
                                   & *+[F]{\begin{array}{l} f=t_{\overrightarrow{u}} \\ \overrightarrow{u} (b) \end{array}}
                                          & 
                                                 & *+[F]{\begin{array}{l} f=r(I;\alpha) \\ \alpha= \mathrm{Arg}\left( a \right) \\ z_I=\dfrac{b}{1-a} \end{array} } \\
	                     }\)
\end{frame}
\begin{frame}{Classification des isométries}
Tout antidéplacement a pour expression complexe une équation de la forme :
                  \(z' = a\overline{z} + b\) avec \(|a| = 1\).

                              Si \(a\overline{b}+b = 0\), \(f\) est une réflexion. 
                        
                              Si \(a\overline{b}+b \neq 0\), il existe une droite \((\mathcal{D})\) et un vecteur 
	                              \(\overrightarrow{u} \neq \overrightarrow{0}\) qui dirige \((\mathcal{D})\) tels que 
                              \(f = t_{\overrightarrow{u}} \circ S_{(\mathcal{D})}\)
                              (forme réduite d'un antidéplacement).
\end{frame}
\begin{frame}{Composition de déplacements et antidéplacements}
On désigne par :

\begin{itemize}

            \item \(s_{(\mathcal{D})}\)  la symétrie orthogonale d'axe \((\mathcal{D})\),
            
\item \(t_{\overrightarrow{u}}\)  la translation de vecteur \(\overrightarrow{u}\),
            
\item \(r(\Omega,\theta)\)  la rotation de centre \(\Omega\) et d'angle \(\theta\).

\end{itemize}
\end{frame}
\begin{frame}{Composition de déplacements et antidéplacements}
Soient \((\mathcal{D}_1)\) et \((\mathcal{D}_2)\) deux droites parallèles distinctes.
            Soient \(H_1\) un point de \((\mathcal{D}_1)\) et \(H_2\) le projeté orthogonal de \(H_1\) sur \((\mathcal{D}_2)\).
            Alors \(s_{(\mathcal{D}_2)} \circ s_{(\mathcal{D}_1)} = t_{2\overrightarrow{H_1 H_2}}\).
\end{frame}
\begin{frame}{Composition de déplacements et antidéplacements (2)}
La réciproque est aussi vraie.
            Soit \(\overrightarrow{u}\) un vecteur et \((\mathcal{D}_1)\) une droite
            perpendiculaire à la direction de \(\overrightarrow{u}\). 
            Soit \((\mathcal{D}_2)\) le translaté de \((\mathcal{D}_1)\) par le vecteur \(\dfrac{1}{2}\overrightarrow{u}\).
            Alors \(t_{\overrightarrow{u}} = s_{(\mathcal{D}_2)} \circ s_{(\mathcal{D}_1)}\).
\end{frame}
\begin{frame}{Composition de déplacements et antidéplacements}
Soient \((\mathcal{D}_1)\) et \((\mathcal{D}_2)\) deux droites distinctes se coupant en un point \(\Omega\).
            Soit \(\theta = \mathrm{mes}((\mathcal{D}_1),(\mathcal{D}_2)) \quad [\pi]\).
            Alors \(s_{(\mathcal{D}_2)} \circ s_{(\mathcal{D}_1)} = r(\Omega,2\theta)\).
\end{frame}
\begin{frame}{Composition de déplacements et antidéplacements (2)}
La réciproque est aussi vraie.
            Soit \(\Omega\) un point et \(\theta\) un nombre réelqui sera considéré en tant qu'angle.
            Soit \((\mathcal{D}_1)\) une droite et soit 
            \((\mathcal{D}_2)\) le transformé de \((\mathcal{D}_1)\) par la rotation de centre \(\Omega\) et d'angle \(\dfrac{\theta}{2}\).
            Alors \(r(\Omega,\theta) = s_{(\mathcal{D}_2)} \circ s_{(\mathcal{D}_1)}\).
\end{frame}
\begin{frame}{Composition de déplacements et antidéplacements}
\textbf{Composée de deux rotations}

 Soient \(r_1 = r(\Omega_1 , \theta_1)\) et \(r_2 = r(\Omega_2 , \theta_2)\).

\begin{itemize}

                        \item Si \(\theta_1 + \theta_2 = 0 \quad [2\pi]\) alors \(r_2 \circ r_1\) est une translation.
                        
\item Si \(\theta_1+\theta_2\neq 0\quad [2\pi]\) alors \(r_2\circ r_1\) est 
                              une rotation d'angle \(\theta_1+\theta_2\).

\end{itemize}
\end{frame}
\begin{frame}{Composition de déplacements et antidéplacements}
\textbf{Composée d'une rotation et d'une translation}

 Soient \(t\) une translation et \(r\) une rotation d'angle \(\theta \neq 0 \quad [2\pi]\).
                  Alors \(r \circ t\) et \(t \circ r\) sont des rotations d'angle \(\theta\) (en général de centres différents).
\end{frame}
\begin{frame}{Composition de déplacements et antidéplacements}
\begin{itemize}

                  \item Tout déplacement est la composée de deux symétries orthogonales.
                  
\item Tout anti-déplacement est soit une symétrie orthogonale, 
                        soit la composée de trois symétries orthogonales.

\end{itemize}

Une symétrie centrale de centre \(\Omega\) est la même chose qu'une
                        rotation de centre \(\Omega\) et d'angle \(\pi\).
\end{frame}
\end{document}
